\chapter{Алгоритм решения задачи}\label{ch:7}

На UML Activity диаграмме показан алгоритм решения задачи (рисунок 17). Система реализует единый сценарий оценки стоимости и длительности ремонтных работ с разветвлением на два режима - автоматический (с использованием моделей компьютерного зрения) и ручной (без использования ML). 

На диаграмме представлены 4 "участника" процесса: Пользователь, Bot Service, Backend и ML Worker. 
Рассмотрим алгоритм взаимодействия пользователя и системы подробнее:

\begin{enumerate}[label=\arabic*.]
    \item Пользователь инициирует взаимодействие с системой, отправляя команду /start или нажимая кнопку "Начать" в боте в мессенджере. Взаимодействие с сервером обеспечивает Bot Service, который передаёт команды из мессенджера системе. На сервере формируется приветственное сообщение и выводится пользователю в мессенджере. Оно содержит клавиатуру с возможностью выбрать режим работы - "С фотографией" и "Ручной ввод". Здесь пользователь выбирает режим, и сценарий разветвляется.
    \item ML-режим: создание заявки. Сценарий переводится в статус CREATED, а Bot Service отображает пользователю сообщение об отправке фотографии повреждённого автомобиля.
    \item После загрузки фотографии она передаётся серверу, валидируется и отправляется в объектное хранилище S3. Если формат был некорректен, пользователю отображается сообщение об ошибке формата и запрашивается повторная отправка изображения. Если в рамках сообщения приложено несколько фотографий (что позволяет сделать ВКонтакте или другой мессенджер), то пользователю отображается сообщение о том, что в рамках системы обрабатывается только первая фотография.
    \item После валидации изображения статус переводится в QUEUED, и формируется задача инференса. Пользователь получает сообщение о начале обработки.
    \item Далее в ML Worker выполняется предобработка и последующий анализ изображения. На первом этапе выполняется детекция и сегментация деталей автомобиля с помощью предобученной модели. Полученные результаты фильтруются по порогу уверенности, и, если деталей не обнаружено - формируется сообщение об ошибке инференса, происходит переход на обработку результатов.
    \item Если детали были обнаружены первой моделью, далее выполняется обрезка этих деталей на фотографии и детекция и сегментация повреждений на каждой из обрезанных областей. После этого формируется аннотированное изображение, и происходит переход на обработку результатов. Если на этапе распознавания деталей или повреждений не получен пригодный результат, система переводит пользователя в сценарий ручного ввода, сохраняя возможность дальнейшего расчета стоимости.
    \item При выборе пользователем ручного режима создаётся заявка со статусом PRICING, минуя этапы инференса с помощью ML-моделей.
    \item Далее происходит переход к фазе, общей для обоих режимов - отображается клавиатура с доступными действиями по формированию / редактированию списка повреждений, в зависимости от выбора пользователь может добавить, удалить, отредактировать повреждение. Отображаются только уникальные сочетания вида <деталь-повреждение>, что упрощает восприятие и редактирование пользователем.
    \item После окончания редактирования пользователь нажимает кнопку "Подтвердить", и система переходит к расчёту стоимости на основании заданных параметров. Статус переводится из PRICING в DONE, а пользователю отображается финальное сообщение о стоимости и длительности работ.
\end{enumerate}

\begin{figure}[h]
\centering
\includegraphics[scale=0.5]{UML_new}
\caption{UML Activity диаграмма сценария системы}
\end{figure}
